\section{Introduction}

Deployed agents rarely run on a bare system prompt. Operators give them
personas (``you are a resourceful engineer who unblocks people''), and more and
more often they work alongside other agents, passing tickets and notes
back and forth \citep{carichon2025multiagent}. Both are social context the
model did not choose. If an agent's willingness to follow a stated rule shifts
with the character it is told to play, or with what other agents say they did,
then a safety evaluation run on the bare model misses how it behaves once
deployed.

Work on conformity shows that LLMs follow a majority on factual and opinion
questions, more strongly for smaller models and confident majorities
\citep{weng2025conformity,zhu2025conformity,choi2025conformity}. Persona
assignment is known to shift toxicity and style \citep{deshpande2023toxicity}
and to act through identifiable directions in activation space
\citep{chen2025persona}. Few studies ask whether either kind of context moves a
\emph{rule-following decision} in an agentic setting, where the behaviour is a
tool call rather than an answer.

That question is easy to measure badly. In an earlier version of this project
(a five-day sprint), we found three ways to do so. (i) A regex checker
misread requests to an approved host as violations. (ii) A capable model ran
code fetches it \emph{expected to be blocked} while refusing in prose. When
those fetches unexpectedly succeeded, it disowned the data 95\% of the time, so
a raw tool-call rate overstated its violations by more than an order of
magnitude. (iii) Most effects rested on three prompts and one model. This paper
is a rebuild designed around those failures:

\begin{itemize}
  \item \textbf{A judge-free say/do measure.} Simulated tools return a
  \emph{successful} payload carrying a unique marker. Whether that marker
  reaches the user separates delivering forbidden data from probing, with a
  string match instead of an LLM judge.
  \item \textbf{A within-design capability control.} Every scenario also runs
  under a \emph{permit} policy, so a low violation rate is read as compliance
  only for models that do make the request when allowed.
  \item \textbf{Persona and peer factors with presence controls}, crossed over
  30 scenarios and \R{all}{nmodels} models from Anthropic, OpenAI and five
  open-weights families.
  \item \textbf{An open Inspect AI implementation} in which every number in
  this paper is generated from the evaluation logs.
\end{itemize}
